\section*{Limitations}

The strict naturally occurring scale-error subsets on official ChartQA are small. We therefore use them as recovery probes and rely on ScaleQA-Natural for broader regime analysis, but ScaleQA-Natural is still a mined diagnostic set rather than a fully independent human-authored benchmark. Its final 520 positive inclusions are author-adjudicated and independently audited by a second human reviewer, but prevalence estimates remain tied to the mining sources, and the source audit shows that the diagnostic items are not an ordinary i.i.d. benchmark split. The main positive adaptation result is obtained with Qwen2.5-VL-7B-Instruct in an answer-tuned setting. We add InternVL2.5-1B, 8B, and 26B diagnostics, including five-seed 8B adaptation, but these experiments show bounded transfer rather than universal improvement. The external PlotQA, ChartBench, MathVista, and InfographicVQA subsets are useful robustness checks, but they are not full-dataset benchmark submissions and should not be interpreted as state-of-the-art comparisons. Our stronger baselines include prompt variants, multi-candidate verification, learned and rule-based selectors, statistical policy fusion, Bayesian bucket/MLP selectors, structured chart-context verification, token-likelihood candidate scoring, and a DePlot-style table-extraction-and-calculation pipeline, but these are still lightweight baselines rather than specialized chart parsers. The scale definition focuses primarily on factor-of-100 style errors, especially percentage-decimal confusions, although ScaleQA-Natural includes smaller groups of 10x, 1000x, and unit-like candidates. Finally, the no-leak selector improves over self-check only modestly and remains far below oracle candidate selection. A deployable system would need stronger scale-regime detection using OCR, axis/unit evidence, visual features, or calibrated verifier probabilities before applying scale correction.
